\section{The comparison with Song and Liang in detail}\label{app:sl}

Section~\ref{sec:assumptions} gives a one-line verdict for each assumption; this appendix gives the argument for the
entries where it takes more than a line. Page numbers are those of \citet{SongLiang2015}; ``supp.\ p.'' is a page of their
online supplement.

\paragraph{\lab{DIM}.} \emph{Mixed.} At the top ours is \emph{stronger}: their Corollary~1 allows $P_n<\exp(Cn^\alpha)$ (p.~954), and their
Theorem~S allows slightly more than polynomial growth (Corollary~S$_1$, $P_n<n^{\kappa\log n}$, supp.\ p.~16). At the
bottom ours is \emph{weaker}: it allows bounded $p$, while their standing choice is $P_n\succ n^\theta$ (p.~950), and
their Theorems~1--3 force $P_n\to\infty$ (their (13), $r_n\prec P_n$ with $r_n\ge1$, pp.~952--953, 956).

\paragraph{\lab{SGR}, \lab{GR}.} \emph{Stronger} than their Condition~1 (p.~953) on centred, unit-variance predictors, in the regime where their
model-size cap $\bar r_n$ is at most our window $5(A+1)(\sz+1)\log n$: ours then implies theirs. Their Condition~1 constrains
combinations of up to $\bar r_n$ predictors, and their Corollary~1 allows $\bar r_n$ beyond any multiple of $\log n$
(p.~954); there neither implies the other, except for centred Gaussian rows with unit variances (their example, p.~953),
which satisfy both. Against their bounded predictors
$\abs{x_i}\le1$ (Theorem~1, p.~952), neither implies the other.

\paragraph{\lab{EFS}.} \emph{Mixed}; it depends on the theorem and on the size of the sets each condition constrains.
\begin{itemize}[leftmargin=1.5em,itemsep=0pt,topsep=1pt]
  \item \emph{Their Condition~3} (pp.~955--956, used by their Theorems~2--3; shown for Gaussian predictors) compares one
    true predictor with combinations of up to $\bar r_n$ others and asks only a vanishing level. Our fixed floor implies it
    when our window covers those sets (the true predictor included); ours is then \emph{stronger}. Otherwise, as when
    $\bar r_n\to\infty$ faster than our window grows, neither implies the other.
  \item \emph{Their Theorem~1} (p.~952) assumes bounded predictors and bounds on the prior covariance, and has no Gram
    floor: not comparable.
  \item \emph{Their Theorem~S} uses (22)--(23) (supp.\ p.~13): (22) bounds the realised Gram matrix both below and above on
    sets of size up to $\min(\sz+\bar r_n,n)$, and (23) is a rate condition involving the lower bound; for random rows they
    relax it to (26) (supp.\ p.~15), with both spectral tails. Ours asks no upper bound, so neither implies the other.
\end{itemize}

\paragraph{Assumption~\ref{asm:dist}.} Items (a)--(c) compare as \as{GL}, \as{DIM}, \as{SP} and \as{SGR}. Item (e) is an
upper bound on each signal in noise units, after scaling each predictor, with no lower bound; it has no direct counterpart
in their theorems (their Theorem~1 controls the coefficients through (3) and (8), pp.~950, 952). Item (d) is
\emph{stronger} than their Theorems~1--2, which have no such condition, and not comparable with their other conditions:
their Condition~3 (p.~955) keeps one true predictor from being imitated by others and puts no floor on sets of null
predictors (two duplicated null columns satisfy it and violate ours), and their (22) (supp.\ p.~13) is a floor, at a higher
level, on the realised design rather than on the population correlation. Our floor, $c_\psi/n$ on the correlation scale, is of the same order as the ridge $c_n\propto1/n$ they use
(p.~953).

\paragraph{Assumption~\ref{asm:reported}.} \emph{None.} \SL{} have no hypothesis of this kind. The nearest text is their informal remark~(a), p.~958 (``if there
are $p_n$ predictors correlated with the true predictor $X_t$, $p_n\prec P_n$, then there are at most $cp_n$ predictors to
be selected''), which is not a condition of any of their theorems; their results on their sampler's output carry no design
condition of this kind (supp.\ pp.~9--10, 12--13, 16--17). The floor~(b) compares as \as{EFS}; the local regressions~(c)
have no counterpart (they serve the method's local search, which \SL{} do not analyse).
